\section{Clausura algebraica del modo de Fourier nulo asociado a la periodicidad CT108 en la fibra electrónica}
\label{sec:clausura-nulo-material-electron-rev3}
\index{electrón!clausura nulo--material}
\index{CT108!modo nulo y electrón}

La identidad autodual de Planck fija el origen de la recta positiva de
masas; la genealogía electrónica fija su primer desplazamiento material
completo. Esta sección reúne, sin repetir sus demostraciones de referencia,
la realización circular de CT108, la composición basal del electrón y la
holonomía bidireccional de su ruta.

Sea
\[
 \mathcal H_{108}^{\rm reg}
 :=\ell^2(\mathbb Z/108\mathbb Z;\mathbb R),
 \qquad
 S_{108}^{\rm reg}\delta_\theta=\delta_{\theta+1},
\]
la representación regular real del desplazamiento cíclico de CT108, y sea
\[
 n_0:=\frac1{\sqrt{108}}\sum_{\theta=0}^{107}\delta_\theta,
 \qquad
 \mathcal N_{108}:=\ker(S_{108}^{\rm reg}-I)=\mathbb R n_0.
\]
Así queda definido, y no sólo nombrado, el modo de Fourier nulo normalizado.
Sea \(\mathcal G_e\) el espacio de genealogías electrónicas enriquecidas de
\cref{prop:planck-ct108-realizacion-centro-electron-rev3} y fijemos el
submódulo lineal
\(\ell_{\Gamma_e}:=\mathbb R g_{\Gamma_e}
\subset\mathbb R[\mathcal G_e]\) generado por la ruta central en la
linealización libre de ese espacio. Asimismo, si
\(F_e^+=\mathbb R_{>0}I\) es la semirrecta positiva ya definida, sea
\(\mathcal F_e:=\operatorname{span}_{\mathbb R}(F_e^+)=\mathbb R I\) su
envolvente lineal, y sea
\(\mathcal L_M\) la recta de masa definida en
\cref{sec:ley-general-masa-accion-autonoma}. Las bases declaradas determinan
el morfismo lineal de rango uno
\begin{equation}
 \mathfrak C_e:
 \mathcal N_{108}\otimes\ell_{\Gamma_e}
 \longrightarrow \mathcal F_e\otimes\mathcal L_M,
 \qquad
 \mathfrak C_e(n_0\otimes g_{\Gamma_e})
 =\mathfrak c_e\,I\otimes\mathcal U_M,
 \label{eq:morfismo-clausura-nulo-material-electron-rev3}
\end{equation}
donde la coordenada positiva es
\begin{equation}
 \boxed{
 \mathfrak c_e=
 \frac{\sqrt3}{4}
 \left[
  \exp\!\left(\frac{\varphi_{\rm HMT}}{\pi_{\rm HMT}^{2}}\right)
  -22\alpha_{\rm HMT}^{3}
 \right]
 (1+15\Delta_4)
 R_{\rm act}^{\,80-54\alpha_{\rm HMT}+6\Delta_4}.}
 \label{eq:coordenada-clausura-nulo-material-electron-rev3}
\end{equation}

\begin{theorem}[Clausura nulo--material electrónica de rango uno]
\label{thm:clausura-nulo-material-electron-rev3}
Con las primitivas y cartas ya fijadas, el morfismo
\(\mathfrak C_e\) factoriza exactamente la salida electrónica beta. Sus
cuatro factores conservan, en este orden, la incompatibilidad suma--producto
de APP; la orientación areal--volumétrica junto con el déficit de frontera;
la torsión y la memoria retenida; y la holonomía bidireccional de la ruta
electrónica. En la carta \(\mathcal U_M=1\,\mathrm{MeV}/c^2\),
\begin{equation}
 \mathfrak C_e(n_0\otimes g_{\Gamma_e})
 =0.5109989506900008086082571648\ldots\,
  I\otimes\mathrm{MeV}/c^2.
 \label{eq:masa-clausura-nulo-material-electron-rev3}
\end{equation}
\end{theorem}

\begin{proof}
El \cref{thm:cierre-electron-holografico-rev6} demuestra que los tres
primeros factores de
\eqref{eq:coordenada-clausura-nulo-material-electron-rev3} componen el
escalar basal positivo \(\mathfrak e_{\rm HMT}\). El
\cref{thm:reduccion-electron-two-way-rev11} obtiene, desde dos lectores
independientes de la ruta,
\[
 K_D(\Gamma_e)=K_\Omega(\Gamma_e)=(80,54,6),
 \qquad
 \kappa_e=80-54\alpha_{\rm HMT}+6\Delta_4.
\]
La razón \(R_{\rm act}>0\) pertenece a la única recta de acción
pre-retorno/post-retorno. Por consiguiente, el cuarto factor es positivo y
la composición coincide con
\eqref{eq:electron-beta-two-way-rev9} y
\eqref{eq:masa-electron-two-way-rev11}. La carta dimensional actúa sólo
después de esta factorización.
\end{proof}

\paragraph{Alcance matemático y realización física.}
El resultado exacto es la realización lineal declarada de la factorización
electrónica sobre el modo nulo, la genealogía y la fibra tipados. El modo
nulo de CT108 no selecciona por sí solo la coordenada \(\mathfrak c_e\):
ésta procede de los cuatro factores demostrados en sus residencias y el
morfismo los reúne en una única clausura algebraica de rango uno. La dinámica electromagnética
ligada conserva como requisitos propios un Hamiltoniano estable, la carga,
el espín \(1/2\), el factor de forma y el acoplamiento gauge; esos objetos no
son premisas de la factorización anterior. La unidad de Planck ocupa el
punto autodual de la recta y el electrón ocupa una sección genealógicamente
desplazada: una interpretación como agujero negro clásico pertenece a otro
tipo de realización y no interviene en este morfismo.
